% Persistas & Partners — annexures.
%
% A notice that demands money is only as good as the proof behind it. The firm
% attaches that proof — receipts, bank statements, WhatsApp screenshots — to the
% notice itself, so the served bundle is one document rather than a letter and a
% pile of loose paper that can be said never to have arrived.
%
% % Persist — shared document preamble.
%
% Every template begins with % Persist — shared document preamble.
%
% Every template begins with % Persist — shared document preamble.
%
% Every template begins with \input{_shared/persist-base}. The PRD lists
% nineteen templates for the Smart Form Compiler; each carrying its own copy of
% the font, colour and geometry setup would mean nineteen places to change when
% the firm's letterhead does, and nineteen chances to miss one.
%
% ENGINE: XeLaTeX. Not pdfLaTeX — see src-tauri/KEEL-RULES.md.
%
% Templates get: A4 geometry, Noto (with Devanagari declared), the Persist
% colour tokens, the L/C/R column types, and \firmfooter. They add only their
% own content.

\usepackage[top=20mm, bottom=20mm, left=22mm, right=22mm]{geometry}

% ── Fonts ────────────────────────────────────────────────────────────────────
% Noto covers ₹ (U+20B9) and Devanagari. pdfLaTeX fails outright on the former.
\usepackage{fontspec}
\setmainfont{Noto Serif}[
  UprightFont = *,
  BoldFont    = * Bold,
  ItalicFont  = * Italic,
]
\setsansfont{Noto Sans}
\setmonofont{Noto Sans Mono}[Scale=MatchLowercase]
% Declared so a Hindi party name or title renders rather than dropping to tofu.
% Devanagari needs its own shaper, not merely a font that has the glyphs.
\newfontfamily\devanagarifont{Noto Serif Devanagari}[Script=Devanagari]
% Wrap Devanagari runs: {\hindi ...}
\newcommand{\hindi}[1]{{\devanagarifont #1}}

% ── Packages ─────────────────────────────────────────────────────────────────
\usepackage{microtype}
\usepackage{booktabs}
\usepackage{tabularx}
\usepackage{array}
\usepackage{longtable}
\usepackage{colortbl}
\usepackage{xcolor}
\usepackage{graphicx}
\usepackage{fancyhdr}
\usepackage{multirow}
\usepackage{hhline}
\usepackage{calc}
\usepackage{enumitem}
\usepackage[hidelinks]{hyperref}
\usepackage{parskip}

% ── Colours — mirror src/design-system/tokens.ts ─────────────────────────────
% A client sees the invoice and the portal side by side; they should look like
% the same firm.
\definecolor{accentBlue}{RGB}{74,101,128}     % --color-accent-primary
\definecolor{terracotta}{RGB}{181,96,74}      % --color-accent-secondary
\definecolor{textPrimary}{RGB}{44,44,42}      % --color-text-primary
\definecolor{textSecondary}{RGB}{107,104,98}  % --color-text-secondary
\definecolor{textTertiary}{RGB}{154,149,144}  % --color-text-tertiary
\definecolor{borderGray}{RGB}{226,221,216}    % --color-border
\definecolor{bgSecondary}{RGB}{240,236,229}   % --color-bg-secondary
\definecolor{statusGreen}{RGB}{74,124,89}     % --color-status-clear
\definecolor{statusUrgent}{RGB}{192,57,43}    % --color-status-urgent

% ── Typography ───────────────────────────────────────────────────────────────
\renewcommand{\familydefault}{\rmdefault}
\setlength{\parindent}{0pt}

% ── Fixed-width column types ─────────────────────────────────────────────────
\newcolumntype{R}[1]{>{\raggedleft\arraybackslash}p{#1}}
\newcolumntype{L}[1]{>{\raggedright\arraybackslash}p{#1}}
\newcolumntype{C}[1]{>{\centering\arraybackslash}p{#1}}

% ── Firm footer ──────────────────────────────────────────────────────────────
% Takes the firm identity as arguments rather than reading {{PLACEHOLDERS}}:
% a shared file that assumed particular placeholder names would force every
% template to declare them, whether or not the document shows them.
%   \firmfooter{name}{gstin}{pan}
\newcommand{\firmfooter}[3]{%
  \pagestyle{fancy}%
  \fancyhf{}%
  \renewcommand{\headrulewidth}{0pt}%
  \fancyfoot[C]{\footnotesize\color{textSecondary}
    Page \thepage\ — #1 — GSTIN: #2 — PAN: #3}%
}

% A footer for documents that are not invoices: no tax identifiers, just the
% firm and the page.
%   \plainfooter{name}
\newcommand{\plainfooter}[1]{%
  \pagestyle{fancy}%
  \fancyhf{}%
  \renewcommand{\headrulewidth}{0pt}%
  \fancyfoot[C]{\footnotesize\color{textSecondary} Page \thepage\ — #1}%
}

% ── Letterhead ───────────────────────────────────────────────────────────────
%   \letterhead{firm name}{address}{contact line}
\newcommand{\letterhead}[3]{%
  {\LARGE\bfseries\color{accentBlue} #1}\par
  \vspace{2pt}
  {\small\color{textSecondary} #2}\par
  {\small\color{textSecondary} #3}\par
  \vspace{2mm}
  {\color{borderGray}\rule{\textwidth}{0.4pt}}
  \vspace{4mm}
}

% ── Section heading in the firm's style ──────────────────────────────────────
\newcommand{\persistsection}[1]{%
  \vspace{3mm}
  {\small\bfseries\color{textSecondary} \MakeUppercase{#1}}\par
  \vspace{1mm}
}
. The PRD lists
% nineteen templates for the Smart Form Compiler; each carrying its own copy of
% the font, colour and geometry setup would mean nineteen places to change when
% the firm's letterhead does, and nineteen chances to miss one.
%
% ENGINE: XeLaTeX. Not pdfLaTeX — see src-tauri/KEEL-RULES.md.
%
% Templates get: A4 geometry, Noto (with Devanagari declared), the Persist
% colour tokens, the L/C/R column types, and \firmfooter. They add only their
% own content.

\usepackage[top=20mm, bottom=20mm, left=22mm, right=22mm]{geometry}

% ── Fonts ────────────────────────────────────────────────────────────────────
% Noto covers ₹ (U+20B9) and Devanagari. pdfLaTeX fails outright on the former.
\usepackage{fontspec}
\setmainfont{Noto Serif}[
  UprightFont = *,
  BoldFont    = * Bold,
  ItalicFont  = * Italic,
]
\setsansfont{Noto Sans}
\setmonofont{Noto Sans Mono}[Scale=MatchLowercase]
% Declared so a Hindi party name or title renders rather than dropping to tofu.
% Devanagari needs its own shaper, not merely a font that has the glyphs.
\newfontfamily\devanagarifont{Noto Serif Devanagari}[Script=Devanagari]
% Wrap Devanagari runs: {\hindi ...}
\newcommand{\hindi}[1]{{\devanagarifont #1}}

% ── Packages ─────────────────────────────────────────────────────────────────
\usepackage{microtype}
\usepackage{booktabs}
\usepackage{tabularx}
\usepackage{array}
\usepackage{longtable}
\usepackage{colortbl}
\usepackage{xcolor}
\usepackage{graphicx}
\usepackage{fancyhdr}
\usepackage{multirow}
\usepackage{hhline}
\usepackage{calc}
\usepackage{enumitem}
\usepackage[hidelinks]{hyperref}
\usepackage{parskip}

% ── Colours — mirror src/design-system/tokens.ts ─────────────────────────────
% A client sees the invoice and the portal side by side; they should look like
% the same firm.
\definecolor{accentBlue}{RGB}{74,101,128}     % --color-accent-primary
\definecolor{terracotta}{RGB}{181,96,74}      % --color-accent-secondary
\definecolor{textPrimary}{RGB}{44,44,42}      % --color-text-primary
\definecolor{textSecondary}{RGB}{107,104,98}  % --color-text-secondary
\definecolor{textTertiary}{RGB}{154,149,144}  % --color-text-tertiary
\definecolor{borderGray}{RGB}{226,221,216}    % --color-border
\definecolor{bgSecondary}{RGB}{240,236,229}   % --color-bg-secondary
\definecolor{statusGreen}{RGB}{74,124,89}     % --color-status-clear
\definecolor{statusUrgent}{RGB}{192,57,43}    % --color-status-urgent

% ── Typography ───────────────────────────────────────────────────────────────
\renewcommand{\familydefault}{\rmdefault}
\setlength{\parindent}{0pt}

% ── Fixed-width column types ─────────────────────────────────────────────────
\newcolumntype{R}[1]{>{\raggedleft\arraybackslash}p{#1}}
\newcolumntype{L}[1]{>{\raggedright\arraybackslash}p{#1}}
\newcolumntype{C}[1]{>{\centering\arraybackslash}p{#1}}

% ── Firm footer ──────────────────────────────────────────────────────────────
% Takes the firm identity as arguments rather than reading {{PLACEHOLDERS}}:
% a shared file that assumed particular placeholder names would force every
% template to declare them, whether or not the document shows them.
%   \firmfooter{name}{gstin}{pan}
\newcommand{\firmfooter}[3]{%
  \pagestyle{fancy}%
  \fancyhf{}%
  \renewcommand{\headrulewidth}{0pt}%
  \fancyfoot[C]{\footnotesize\color{textSecondary}
    Page \thepage\ — #1 — GSTIN: #2 — PAN: #3}%
}

% A footer for documents that are not invoices: no tax identifiers, just the
% firm and the page.
%   \plainfooter{name}
\newcommand{\plainfooter}[1]{%
  \pagestyle{fancy}%
  \fancyhf{}%
  \renewcommand{\headrulewidth}{0pt}%
  \fancyfoot[C]{\footnotesize\color{textSecondary} Page \thepage\ — #1}%
}

% ── Letterhead ───────────────────────────────────────────────────────────────
%   \letterhead{firm name}{address}{contact line}
\newcommand{\letterhead}[3]{%
  {\LARGE\bfseries\color{accentBlue} #1}\par
  \vspace{2pt}
  {\small\color{textSecondary} #2}\par
  {\small\color{textSecondary} #3}\par
  \vspace{2mm}
  {\color{borderGray}\rule{\textwidth}{0.4pt}}
  \vspace{4mm}
}

% ── Section heading in the firm's style ──────────────────────────────────────
\newcommand{\persistsection}[1]{%
  \vspace{3mm}
  {\small\bfseries\color{textSecondary} \MakeUppercase{#1}}\par
  \vspace{1mm}
}
. The PRD lists
% nineteen templates for the Smart Form Compiler; each carrying its own copy of
% the font, colour and geometry setup would mean nineteen places to change when
% the firm's letterhead does, and nineteen chances to miss one.
%
% ENGINE: XeLaTeX. Not pdfLaTeX — see src-tauri/KEEL-RULES.md.
%
% Templates get: A4 geometry, Noto (with Devanagari declared), the Persist
% colour tokens, the L/C/R column types, and \firmfooter. They add only their
% own content.

\usepackage[top=20mm, bottom=20mm, left=22mm, right=22mm]{geometry}

% ── Fonts ────────────────────────────────────────────────────────────────────
% Noto covers ₹ (U+20B9) and Devanagari. pdfLaTeX fails outright on the former.
\usepackage{fontspec}
\setmainfont{Noto Serif}[
  UprightFont = *,
  BoldFont    = * Bold,
  ItalicFont  = * Italic,
]
\setsansfont{Noto Sans}
\setmonofont{Noto Sans Mono}[Scale=MatchLowercase]
% Declared so a Hindi party name or title renders rather than dropping to tofu.
% Devanagari needs its own shaper, not merely a font that has the glyphs.
\newfontfamily\devanagarifont{Noto Serif Devanagari}[Script=Devanagari]
% Wrap Devanagari runs: {\hindi ...}
\newcommand{\hindi}[1]{{\devanagarifont #1}}

% ── Packages ─────────────────────────────────────────────────────────────────
\usepackage{microtype}
\usepackage{booktabs}
\usepackage{tabularx}
\usepackage{array}
\usepackage{longtable}
\usepackage{colortbl}
\usepackage{xcolor}
\usepackage{graphicx}
\usepackage{fancyhdr}
\usepackage{multirow}
\usepackage{hhline}
\usepackage{calc}
\usepackage{enumitem}
\usepackage[hidelinks]{hyperref}
\usepackage{parskip}

% ── Colours — mirror src/design-system/tokens.ts ─────────────────────────────
% A client sees the invoice and the portal side by side; they should look like
% the same firm.
\definecolor{accentBlue}{RGB}{74,101,128}     % --color-accent-primary
\definecolor{terracotta}{RGB}{181,96,74}      % --color-accent-secondary
\definecolor{textPrimary}{RGB}{44,44,42}      % --color-text-primary
\definecolor{textSecondary}{RGB}{107,104,98}  % --color-text-secondary
\definecolor{textTertiary}{RGB}{154,149,144}  % --color-text-tertiary
\definecolor{borderGray}{RGB}{226,221,216}    % --color-border
\definecolor{bgSecondary}{RGB}{240,236,229}   % --color-bg-secondary
\definecolor{statusGreen}{RGB}{74,124,89}     % --color-status-clear
\definecolor{statusUrgent}{RGB}{192,57,43}    % --color-status-urgent

% ── Typography ───────────────────────────────────────────────────────────────
\renewcommand{\familydefault}{\rmdefault}
\setlength{\parindent}{0pt}

% ── Fixed-width column types ─────────────────────────────────────────────────
\newcolumntype{R}[1]{>{\raggedleft\arraybackslash}p{#1}}
\newcolumntype{L}[1]{>{\raggedright\arraybackslash}p{#1}}
\newcolumntype{C}[1]{>{\centering\arraybackslash}p{#1}}

% ── Firm footer ──────────────────────────────────────────────────────────────
% Takes the firm identity as arguments rather than reading {{PLACEHOLDERS}}:
% a shared file that assumed particular placeholder names would force every
% template to declare them, whether or not the document shows them.
%   \firmfooter{name}{gstin}{pan}
\newcommand{\firmfooter}[3]{%
  \pagestyle{fancy}%
  \fancyhf{}%
  \renewcommand{\headrulewidth}{0pt}%
  \fancyfoot[C]{\footnotesize\color{textSecondary}
    Page \thepage\ — #1 — GSTIN: #2 — PAN: #3}%
}

% A footer for documents that are not invoices: no tax identifiers, just the
% firm and the page.
%   \plainfooter{name}
\newcommand{\plainfooter}[1]{%
  \pagestyle{fancy}%
  \fancyhf{}%
  \renewcommand{\headrulewidth}{0pt}%
  \fancyfoot[C]{\footnotesize\color{textSecondary} Page \thepage\ — #1}%
}

% ── Letterhead ───────────────────────────────────────────────────────────────
%   \letterhead{firm name}{address}{contact line}
\newcommand{\letterhead}[3]{%
  {\LARGE\bfseries\color{accentBlue} #1}\par
  \vspace{2pt}
  {\small\color{textSecondary} #2}\par
  {\small\color{textSecondary} #3}\par
  \vspace{2mm}
  {\color{borderGray}\rule{\textwidth}{0.4pt}}
  \vspace{4mm}
}

% ── Section heading in the firm's style ──────────────────────────────────────
\newcommand{\persistsection}[1]{%
  \vspace{3mm}
  {\small\bfseries\color{textSecondary} \MakeUppercase{#1}}\par
  \vspace{1mm}
}
 first; this uses its colours and fonts.
%
% MARKING
%
% Every annexure page carries its mark in the top-right corner. That marking is
% not the attorney's job: Keel allocates the marks in order and builds both the
% list printed in the notice and the pages appended after it from the same
% source, so the list cannot say Annexure C while the pages stop at B.
%
% The mark is drawn with eso-pic's starred \AddToShipoutPictureFG*, which is
% "this page only" and expands its argument where it is called, rather than with
% a fancyhdr page style — which is the obvious way, and puts the mark in the
% wrong place. \applyletterhead gives the document a 108pt head band so the
% firm's mark and both partners fit; an annexure page inheriting that geometry
% has its header pushed roughly 40% of the way down the paper, landing on top of
% the first line of the receipt it is supposed to be marking. Compiled and
% measured both ways: fancyhdr put the mark at y=137pt with the receipt's own
% text at y=144pt. A shipout picture is positioned against the paper, not
% against a text area the letterhead controls, so it sits 13mm from the edge
% whatever the body geometry is.
%
% `the_mark_sits_at_the_top_of_the_page_clear_of_the_content` in
% services/annexures.rs is the test that fails if this is changed back.
%
% The annexure pages themselves are \thispagestyle{empty}: the firm's letterhead
% belongs on the firm's words, not overprinted on a third party's receipt.

\usepackage{pdfpages}
\usepackage{eso-pic}

% ── The mark ─────────────────────────────────────────────────────────────────
% Boxed, top-right, 13mm down from the paper edge — clear of the content of a
% scanned A4 page, which almost always has a margin of its own.
%
%   \annexurestamp{A(i)}
\newcommand{\annexurestamp}[1]{%
  \AddToShipoutPictureFG*{%
    \AtPageUpperLeft{%
      \raisebox{-13mm}[0pt][0pt]{%
        \makebox[\paperwidth][r]{%
          {\sffamily\bfseries\fontsize{9}{11}\selectfont\color{textPrimary}%
           \setlength{\fboxsep}{4pt}%
           \setlength{\fboxrule}{0.4pt}%
           \fcolorbox{textPrimary}{white}{\strut ANNEXURE #1}}%
          \hspace{14mm}}}}%
  }%
}

% ── A PDF annexure ───────────────────────────────────────────────────────────
% Every page of the source, scaled to leave room for the mark. The source keeps
% its own page numbering; renumbering someone else's document would misrepresent
% it.
%
%   \annexurepdf{A}{annexure-01.pdf}
\newcommand{\annexurepdf}[2]{%
  \includepdf[pages=-, scale=0.86, offset=0 -8mm,
              pagecommand={\annexurestamp{#1}\thispagestyle{empty}}]{#2}%
}

% ── An image annexure ────────────────────────────────────────────────────────
% A screenshot or a photographed receipt, filling its page. Nothing is printed
% alongside it — not a caption, not a heading. An annexure page carries the
% document and its mark and nothing else, so what the recipient sees is the
% evidence as it was, with only the firm's marking added.
%
%   \annexureimage{B}{annexure-02.png}
\newcommand{\annexureimage}[2]{%
  \clearpage
  \annexurestamp{#1}%
  \thispagestyle{empty}%
  \null\vfill
  \begin{center}
    \includegraphics[width=\textwidth, height=0.86\textheight, keepaspectratio]{#2}
  \end{center}
  \vfill\null
  \clearpage
}

% ── The list printed in the notice ───────────────────────────────────────────
% Set on the last page of the notice body, above the annexures themselves.
% \annexurelistheading is emitted by Keel only when there is at least one
% annexure — a notice with nothing attached should not print an empty heading.
% The \par is load-bearing: without it the first \annexureentry continues the
% heading's paragraph and prints on the same line as "Enclosures / Annexures:".
\newcommand{\annexurelistheading}{%
  \vspace{8mm}
  \noindent\textbf{Enclosures / Annexures:}\par
  \vspace{3mm}
}

%   \annexureentry{A}{Receipt dated 14 July 2026}
\newcommand{\annexureentry}[2]{%
  \noindent
  \makebox[32mm][l]{\textbf{Annexure-#1}}%
  \begin{minipage}[t]{\dimexpr\linewidth-32mm\relax}
    #2
  \end{minipage}\par
  \vspace{1.5mm}
}
